% 360_HRV_Galois_Experiment_De_ch.tex
% Johann Pascher, 8. September 2026
% Galois-informierte HRV-Analyse: Implementierung und explorative Bewertung
\providecommand{\BM}{\textbf{[B]}}
\providecommand{\KM}{\textbf{[K]}}
\providecommand{\SM}{\textbf{[S]}}
\providecommand{\XM}{\textbf{[X]}}
\providecommand{\LM}{\textbf{[L]}}

\phantomsection
\addcontentsline{toc}{chapter}{%
  Galois-informierte HRV-Analyse --- Implementierung und explorative Bewertung}

\section*{Statuszeichen}
\begin{center}
\begin{tabular}{@{}lp{10.5cm}@{}}
\textbf{[B]} & \emph{Belegt} --- algebraisch bewiesen. \\[3pt]
\textbf{[K]} & \emph{Konfirmiert} --- numerisch geprüft. \\[3pt]
\textbf{[S]} & \emph{Spekulativ} --- offen. \\[3pt]
\textbf{[L]} & \emph{Literaturbeleg} --- externer Primärnachweis. \\
\end{tabular}
\end{center}

\section*{Abstract}

Dok.~358 und 359 formulieren eine algebraische Hypothese für die
HRV-Frequenzanalyse: Frequenzverhältnisse gekoppelter physiologischer
Oszillatoren sollen auf einem Galois-Raster liegen, mit der
Strukturtoleranz $100\xi\approx1{,}33\,\%$. Dieses Dokument setzt den
Ansatz als Python-Modul um und bewertet ihn explorativ an zwei
öffentlichen EKG-Datensätzen und an eigenen Polar-H10-Aufnahmen. Das
Ergebnis ist ein Negativbefund: Die vorliegenden Daten stützen die
Hypothese nicht, und sie können sie auch nicht prüfen, weil LF/HF ein
Leistungs- und kein Frequenzverhältnis ist.

\textbf{Ergebnisse:}
\begin{enumerate}
\item \textbf{Modul} \KM: \texttt{galois\_hrv\_preprocess.py} ---
  Galois-Projektion, HRV-Bandgrenzen, Orbit-Tiefe, Komplexitätsmaß
  $2/(p+q)$.
\item \textbf{Nullmodell} \KM: Im Raster bis Nenner 60 liegt jeder Wert
  zwischen 1 und 4 innerhalb von 1{,}33\,\% eines Rasterelements;
  innerhalb von 0{,}06\,\% liegen 39\,\% zufälliger Werte. Nahe
  Gittertreffer sind daher für sich kein Beleg.
\item \textbf{BIDMC-01} \LM: Intensivpatient, 88~J., 8~min. Kein
  Verhältnis erfüllt das 5-Limit-Kriterium; der Datensatz ist für die
  Frage ungeeignet.
\item \textbf{Fantasia} \LM: 5 junge und 5 alte Gesunde, je
  $\approx120$~min, spontane Atmung. LF-Peak-CV 11--33\,\%; 5-Limit-LF/HF
  4/10 bei einer Nullrate von 21\,\% (Binomialtest $p=0{,}14$).
\item \textbf{Polar H10} \KM: Mit validierter Schlagerkennung liegt der
  dominante Spektralpeak bei geführter Atmung (4--6/min) innerhalb einer
  Frequenzstufe auf der Atemfrequenz, der HF-Peak ist ihre zweite
  Harmonische. Gittertreffer von LF/HF treten nur mit Zufallshäufigkeit auf.
\item \textbf{Struktureller Grund} \BM: Bei spontaner Atmung ist der
  HF-Peak die Atemfrequenz und LF/HF das Verhältnis zweier unabhängiger
  Regelkreise; bei geführter Atmung unter 9/min liegt die Atmung selbst
  im LF-Band.
\item \textbf{Nächster Schritt} \SM: prospektiver Test der
  Frequenzkopplung bei geführter Atmung (Abschnitt~\ref{sec:360next}).
\end{enumerate}

% ============================================================
\section{Galois-HRV-Preprocessing-Modul}
% ============================================================

\texttt{galois\_hrv\_preprocess.py} implementiert (Dok.~359 §5):

\begin{description}
\item[Galois-Raster] 901 rationale Verhältnisse $p/q\le4$ mit Primfaktoren
  $\subseteq\{2,3,5,7,11,13\}$ bis Nenner 60. Die Primzahlmenge ist
  $\{3\}\cup\{p : p\mid 3^k-1,\ k\le6\}$; eine Primzahl $p\neq3$ gehört zur
  Tiefe $k$, wenn $k$ ein Vielfaches von $\mathrm{ord}_p(3)$ ist.
\item[Toleranz] $100\xi = 1/75 \approx 1{,}33\,\%$ mit Farey-Höhe
  $Q_c=\pi/\sqrt{6\cdot100\xi}\approx11{,}1$ (Dok.~343): Bei diesem Nenner
  ist der mittlere Farey-Abstand gleich der doppelten Toleranz. Ein Raster
  bis Nenner 60 ist dichter als die Toleranz und trennt nicht.
\item[HRV-Bandgrenzen] VLF: $[0,\,\tfrac{1}{25}]$~Hz, LF: $[\tfrac{1}{25},\,\tfrac{3}{20}]$~Hz,
  HF: $[\tfrac{3}{20},\,\tfrac{2}{5}]$~Hz. Die Primfaktoren liegen in
  $\{2,3,5\}$ \BM; die Grenzen sind damit mit dem 5-Limit-Raster verträglich,
  aber nicht durch es festgelegt (Abschnitt~\ref{sec:360baender}).
\item[Galois-Projektion] Verhältnis $\mapsto$ nächstes $p/q$ im Raster,
  Abstand $d=|r/(p/q)-1|$.
\item[Orbit-Tiefe] $k=\max\,\mathrm{ord}_p(3)$ über die Primfaktoren
  $p\notin\{2,3\}$ (Dok.~343, 358).
\item[Komplexitätsmaß] $2/(p+q)$ ordnet Verhältnisse nach ihrer
  Einfachheit (Dok.~328). Es ist keine berechnete Arnold-Zungenbreite;
  dazu wäre die Kopplungsstärke eines konkreten Oszillatormodells nötig.
\item[Pipeline] \texttt{hrv\_galois\_pipeline(rr\_ms)} --- ein Funktionsaufruf.
\end{description}

\textbf{Auflösungsbedingung.} Eine Klassifikation mit der Toleranz
$\delta=100\xi$ ist nur sinnvoll, wenn die Frequenzauflösung
$\Delta f/f\le\delta$ erfüllt. Mit $\Delta f=1/T_w$ für die Fensterlänge
$T_w$ (bei Welch die Segmentlänge, nicht die Aufnahmedauer) folgt
$T_w\ge1/(\delta f)$: 12{,}5~min bei 0{,}10~Hz, 15~min bei 0{,}083~Hz \BM.
Davon zu unterscheiden sind der Mindestabstand zweier trennbarer Peaks
(etwa $2/T_w$ beim Hann-Fenster) und die Genauigkeit der Lage eines
einzelnen Peaks, die durch Interpolation feiner als $\Delta f$ werden kann.

\textbf{Frequenz- und Leistungsverhältnisse.} Die Hypothese betrifft
Frequenzverhältnisse gekoppelter Oszillatoren. LF/HF ist ein Verhältnis
integrierter Bandleistungen; es gibt keinen Grund, es auf dem Raster zu
erwarten. Die Projektion von LF/HF dient im Folgenden nur dazu, das
Verhalten des Nullmodells zu zeigen.

% ============================================================
\section{Bandgrenzen und Nullmodell}\label{sec:360baender}
% ============================================================

\textbf{Bandgrenzen} \KM. Die drei Task-Force-Grenzen haben Primfaktoren in
$\{2,3,5\}$ in Hz, mHz, Schlägen pro Minute ($2{,}4=12/5$, 9, 24) und als
Perioden (25~s, $20/3$~s, $5/2$~s), weil die Umrechnungsfaktoren 60 und 1000
nur diese Primzahlen enthalten. In rad/s sind sie irrational und nicht
klassifizierbar. Die Eigenschaft ist auch nicht spezifisch: 8 von 9 Zahlen
mit einer gültigen Ziffer und ein Drittel aller zweistelligen Dezimalwerte
in $(0,1)$ sind 5-Limit. Gerundete Konventionen landen damit von selbst im
Raster; die Bandgrenzen sind verträglich, aber nicht algebraisch festgelegt.
Frequenzverhältnisse sind dagegen dimensionslos und einheitenunabhängig.

\textbf{Nullmodell} \KM. Wahrscheinlichkeit, dass ein log-gleichverteilter
Wert in $[1,4]$ innerhalb der Toleranz $\tau$ eines Rasterelements liegt
(200\,000 Ziehungen):

\begin{center}
\begin{tabular}{@{}lrr@{}}
\toprule
$\tau$ & Raster bis Nenner 60 (901) & Raster bis Nenner 11 (133) \\
\midrule
1{,}33\,\% & 100{,}0\,\% & 90{,}5\,\% \\
0{,}14\,\% & 73{,}9\,\% & 18{,}5\,\% \\
0{,}06\,\% & 39{,}5\,\% & 7{,}9\,\% \\
0{,}02\,\% & 14{,}5\,\% & 2{,}6\,\% \\
\bottomrule
\end{tabular}
\end{center}

Mit dem Raster bis Nenner 60 kann das Kriterium \glqq innerhalb 1{,}33\,\%\grqq{} nichts
unterscheiden. Auch das Raster bis zur Farey-Höhe 11 deckt bei 1{,}33\,\%
noch 91\,\% ab.

% ============================================================
\section{Negativbefund I: BIDMC-01}
% ============================================================

\textbf{Datensatz:} BIDMC PPG and Respiration Dataset v1.0.0,
Record 01 \LM\,[Pim18].
88~J., m., MICU, 8~min, 125~Hz PPG, EKG und Respiration.

\textbf{Ergebnis:} 0/3 Frequenzverhältnisse erfüllen das 5-Limit-Kriterium
(Welch und Lomb-Scargle). VLF ist in 8~min nicht auflösbar. DFA
$\alpha_1=0{,}50$, $\alpha_2=0{,}45$ --- weißes Rauschen; bei Gesunden
$\approx1{,}0$.

\textbf{Warum:}
\begin{enumerate}
\item Zu kurz für VLF.
\item Intensivpatient: fraktale Regulation weitgehend verloren.
\item HF-Peak = gemessene Atemfrequenz (0{,}328~Hz gegenüber 0{,}336~Hz im
  RESP-Kanal).
\item Messunsicherheit Welch gegenüber Lomb-Scargle $\pm2\,\%$, größer als
  die Toleranz 1{,}33\,\%.
\end{enumerate}

% ============================================================
\section{Negativbefund II: Fantasia}
% ============================================================

\textbf{Datensatz:} PhysioNet Fantasia Database v1.0.0
\LM\,[Iye96]. 5 jung (f1y01--05, 21--34~J.) + 5 alt
(f2o01--05, 68--85~J.), je $\approx120$~min Ruhe, 250~Hz EKG,
spontane Atmung. Referenz-Schlagannotationen (.ecg), CC-BY 4.0.

\textbf{Methodik:} wfdb-Annotationen, Welch 10-min-Fenster,
Zero-Padding $\times4$, parabolische Peak-Interpolation,
LF-Stabilität über 6 $\times$ 5-min-Segmente.

\begin{center}
\begin{tabular}{@{}lrrrrrrrl@{}}
\toprule
Rec & $n_\text{RR}$ & LF~Hz & CV\,\% & LF/HF & $p/q$ & $k$ & 5-Lim & $2/(p+q)$ \\
\midrule
f1y01 & 8707 & 0{,}109 & 11{,}5 & 0{,}438 & $7/16$ & 6 & N & 0{,}087 \\
f1y02 & 6919 & 0{,}085 & 33{,}0 & 0{,}867 & $13/15$ & 4 & N & 0{,}071 \\
f1y03 & 7641 & 0{,}100 & 13{,}7 & 4{,}219 & $4$ & 1 & J & 0{,}400 \\
f1y04 & 5232 & 0{,}078 & 30{,}9 & 1{,}060 & $35/33$ & 6 & N & 0{,}029 \\
f1y05 & 6943 & 0{,}068 & 16{,}5 & 2{,}993 & $3$ & 1 & J & 0{,}500 \\
\midrule
f2o01 & 7045 & 0{,}067 & 25{,}6 & 1{,}494 & $3/2$ & 1 & J & 0{,}400 \\
f2o02 & 6327 & 0{,}052 & 14{,}4 & 0{,}864 & $45/52$ & 4 & N & 0{,}021 \\
f2o03 & 6478 & 0{,}055 & 22{,}0 & 0{,}841 & $21/25$ & 6 & N & 0{,}043 \\
f2o04 & 6849 & 0{,}065 & 32{,}2 & 0{,}719 & $28/39$ & 6 & N & 0{,}030 \\
f2o05 & 8334 & 0{,}049 & 11{,}4 & 1{,}873 & $15/8$ & 4 & J & 0{,}087 \\
\bottomrule
\end{tabular}
\end{center}

\textbf{Statistische Tests (Mann-Whitney~U):} Alle Gruppenunterschiede
jung gegenüber alt n.s.\ ($p>0{,}40$).

\textbf{Befunde:}
\begin{itemize}
\item LF-Peak-Stabilität: 0/10 Records stabil ($\text{CV}>1{,}33\,\%$).
\item 5-Limit-LF/HF (nächstes Rasterelement hat Primfaktoren in $\{2,3,5\}$):
  4/10. Die Nullrate für dieses Kriterium ist 21\,\%; einseitiger
  Binomialtest $p=0{,}14$, also nicht über Zufall.
\item Kein Gruppenunterschied jung/alt ($n=5$ je Gruppe).
\end{itemize}

% ============================================================
\section{Polar H10: Aufnahmen und Auswertung}
% ============================================================

\textbf{Gerät und Protokoll:} Polar H10 EKG-Brustgurt, Firmware 4.2.0,
JSONL-Export über Polar Sensor Recording (Android), Autor als einziger
Proband (75~J.), 10.~September 2026. Eine Spontanaufnahme liegend
(21{,}9~min, im Gerät berechnete RR-Intervalle), sechs Aufnahmen sitzend
mit Metronom, je 5{,}0--6{,}0~min, 130~Hz EKG. Jeder Klick markiert einen
Wechsel der Atemphase; 4, 8, 10 und 12 Klicks/min entsprechen 2, 4, 5 und
6 Atemzügen/min (0{,}033 bis 0{,}100~Hz). Die beiden als 4/min
bezeichneten Dateien sind byte-identisch; es gibt eine Aufnahme mit
4/min. Die Atmung wurde nicht mit einem eigenen Sensor gemessen.

\textbf{Auswertung} \KM: R-Zacken mit NeuroKit2 \LM\,[Mak21],
Artefaktkriterium $|\text{RR}/\text{Median}-1|>0{,}30$; markierte
Intervalle werden entfernt, die übrigen Schläge bleiben auf ihrer echten
Zeitmarke. Kubische Interpolation auf 4~Hz, lineare Trendbereinigung,
Hann-Periodogramm über die ganze Aufnahme ($\Delta f=1/T\approx0{,}003$~Hz),
Bandleistungen durch Trapezintegration über LF $[0{,}04,\,0{,}15)$ und
HF $[0{,}15,\,0{,}40)$~Hz.

\textbf{Schlagerkennung und Zeitachse.} Der einfache Energie-Detektor in
\texttt{polar\_h10\_atemfrequenz\_scan.py} (fester Schwellwert beim
90.~Perzentil) fand je Aufnahme 216--285 Schläge, der validierte Detektor
396--443. In jeder Aufnahme mit Atemvorgabe fehlten damit 35--47\,\% der
Schläge. Werden die verbleibenden Intervalle aneinandergehängt, schrumpft
die Zeitachse auf 49--61\,\% der Aufnahmedauer, und alle Frequenzen
verschieben sich um den Faktor 1{,}6--2{,}0 nach oben. Die Tabelle beruht
auf der validierten Erkennung mit erhaltener Zeitachse.

\begin{center}
\small
\begin{tabular}{@{}lrrrrrrrl@{}}
\toprule
Bedingung & $f_R$~Hz & $T$~min & Schläge & RMSSD~ms & LF-Peak & HF-Peak & LF/HF & nächstes $p/q$ \\
\midrule
Spontan, liegend & --- & 21{,}9 & 1517~RR & 19{,}0 & 0{,}074 & 0{,}225 & 1{,}65 & $33/20$ (0{,}09\,\%) \\
2/min, Aufn.\,1 & 0{,}033 & 5{,}4 & 431 & 34{,}2 & 0{,}067 & 0{,}200 & 7{,}44 & --- \\
2/min, Aufn.\,2 & 0{,}033 & 5{,}1 & 411 & 49{,}6 & 0{,}067 & 0{,}167 & 4{,}24 & --- \\
4/min & 0{,}067 & 5{,}3 & 396 & 67{,}0 & 0{,}067 & 0{,}358 & 4{,}61 & --- \\
5/min, Aufn.\,1 & 0{,}083 & 5{,}6 & 431 & 64{,}6 & 0{,}083 & 0{,}167 & 3{,}28 & $105/32$ (0{,}006\,\%) \\
5/min, Aufn.\,2 & 0{,}083 & 5{,}1 & 400 & 68{,}5 & 0{,}083 & 0{,}167 & 2{,}47 & $121/49$ (0{,}05\,\%) \\
6/min & 0{,}100 & 6{,}0 & 443 & 60{,}9 & 0{,}100 & 0{,}201 & 4{,}20 & --- \\
\bottomrule
\end{tabular}
\end{center}

\noindent{\small Frequenzen in Hz. ---: LF/HF über 4 (außerhalb des Rasters)
oder mehr als 1{,}33\,\% vom nächsten Element entfernt. Artefakte: je 0--1.}

\textbf{Befund 1 --- Peak auf der Atemfrequenz} \KM: Bei 4, 5 und 6/min
liegt der dominante Peak innerhalb einer Frequenzstufe auf der
Metronomfrequenz. Die Atmung liegt damit im LF-Band: Der LF-Peak
\emph{ist} der Atempeak. Wo die zweite Harmonische über 0{,}15~Hz liegt
(5 und 6/min), ist sie der HF-Peak.

\textbf{Befund 2 --- Resonanzatmung} \KM: RMSSD liegt bei 4--6/min bei
61--69~ms, bei 2/min bei 34--50~ms. Das entspricht der bekannten
Verstärkung der HRV durch langsame Atmung nahe der Barorezeptor-Resonanz
\LM\,[Leh14, Sev22]. Die Spontanaufnahme ist liegend
entstanden und mit den sitzenden Aufnahmen nicht direkt vergleichbar.

\textbf{Befund 3 --- 2/min} \SM: Der dominante Peak liegt bei 0{,}067~Hz,
dem Doppelten der vorgesehenen Atemfrequenz, mit einer kleineren
Komponente bei 0{,}033~Hz. Ohne eigenes Atemsignal ist nicht zu
entscheiden, ob das die zweite Harmonische sehr langsamer Atmung ist oder
Atmung mit 4/min wegen der Klick-Konvention.

\textbf{Befund 4 --- Gittertreffer im Zufallsbereich} \KM: Die beiden
Aufnahmen bei 5/min liegen nahe an zwei \emph{verschiedenen}
Rasterelementen ($105/32$ und $121/49$), obwohl die Bedingung dieselbe
ist; ihr LF/HF unterscheidet sich um 28\,\%. In 17 gleitenden
300-s-Fenstern der Spontanaufnahme schwankt LF/HF zwischen 0{,}89 und
7{,}43 (CV 78\,\%); von den 13 Fenstern im Rasterbereich $[1,4]$ liegen 6
innerhalb 0{,}06\,\% eines Elements, bei einer Nullrate von 39\,\%
(Binomialtest $p=0{,}78$). Die Streuung der Schätzung ist um mehr als drei
Größenordnungen größer als die Abstände zum Raster.

\textbf{Spontanaufnahme} \KM: mittlere Herzfrequenz 69~/min, RMSSD
19{,}0~ms, SDNN 48{,}9~ms. Ein einzelner verpasster Schlag (1741~ms,
doppelter Median) hebt RMSSD ohne Artefaktbehandlung auf 36{,}3~ms. Der
LF-Peak schwankt über die 300-s-Fenster zwischen 0{,}040 und 0{,}085~Hz
(CV 22\,\%), weit über 1{,}33\,\%. LF/HF beträgt 1{,}65 mit einem Fenster
über die ganze Aufnahme und 2{,}21 mit 300-s-Segmenten; schon eine
Einstellung des Schätzers ändert den Wert um ein Drittel.

\textbf{Methodische Lehren:}
\begin{enumerate}
\item Schlagerkennung gegen die erwartete Schlagzahl
  (Herzfrequenz $\times$ Dauer) und einen validierten Detektor prüfen.
\item Zeitachse erhalten: verworfene Intervalle entfernen, die übrigen
  Schläge auf ihrer Zeitmarke lassen, nie aneinanderhängen.
\item Ein einzelnes Artefakt kann RMSSD dominieren; Artefaktbehandlung ist
  bei jeder Aufnahmedauer angebracht.
\item Unter 9 Atemzügen/min liegt die Atmung im LF-Band; LF/HF hat dann
  keine Deutung als sympathovagale Balance.
\item $n=1$ und 5--6~min Aufnahmen erlauben keine statistischen Schlüsse;
  alle Polar-Befunde sind explorativ.
\end{enumerate}

% ============================================================
\section{Struktureller Grund und Konsequenz}
% ============================================================

\begin{quote}
\textbf{Befund.} \BM\ Bei spontaner Atmung ist der HF-Peak die
Atemfrequenz (variabel 0{,}15--0{,}40~Hz). LF/HF ist das Verhältnis
zweier unabhängiger physiologischer Regelkreise (Barorezeptor~/ Atmung)
und konvergiert nicht gegen ein festes Verhältnis. Bei geführter Atmung
gibt das Metronom die Frequenz vor; Peak und Harmonische stehen dann in
trivialen ganzzahligen Verhältnissen. In beiden Fällen ist LF/HF kein
Prüfgegenstand der Galois-Hypothese.
\end{quote}

% ============================================================
\section{Nächster Schritt: prospektiver Test}\label{sec:360next}
% ============================================================

Die Hypothese sagt etwas über Frequenzkopplung aus. Prüfbar ist das,
wenn die Atemfrequenz gezielt relativ zur individuellen
Barorezeptor-Resonanz $f_\text{res}$ gesetzt wird \SM:

\begin{itemize}
\item Resonanzbestimmung: geführte Atmung 4{,}5--7/min in Schritten von
  0{,}5/min, je 2~min \LM\,[Leh14].
\item Vier Blöcke zu 21~min (Auflösungsbedingung bei der niedrigsten
  Atemfrequenz von etwa 0{,}062~Hz) bei $\rho=f_R/f_\text{res}\in\{2/3,\,3/2\}$
  (Raster) und $\{0{,}618,\,1{,}618\}$ (goldener Schnitt), nach Atemrate
  paarweise abgestimmt, in zufälliger Reihenfolge.
\item Primärer Endpunkt: $n{:}m$-Phasensynchronisationsindex
  \LM\,[Tas98] zwischen Atemphase (Atemgurt) und HRV-Komponente bei
  $f_\text{res}$; Kontrast Raster minus goldener Schnitt, gepaarter
  zweiseitiger $t$-Test.
\item Fallzahl: $d_z=0{,}6$, Power 0{,}80 $\Rightarrow n=24$; mit
  20\,\% Ausfall 30 Teilnehmende.
\item Falsifikation: Ist die Kopplung bei Rasterverhältnissen nicht
  stärker als beim goldenen Schnitt, hat das Komplexitätsmaß $2/(p+q)$ in
  der HRV keine Grundlage.
\end{itemize}

Ein positives Ergebnis stützt die Prämisse, unterscheidet aber noch nicht
die Primzahlmenge und die Farey-Höhe der FFGFT von allgemeinem
Mode-Locking gekoppelter Oszillatoren; dafür dienen sekundäre Endpunkte
(gemessenes Frequenzverhältnis gegen das Raster bis Nenner 11,
Peak-Stabilität gegen 1{,}33\,\%).

% ============================================================
\section{Zusammenfassung}
% ============================================================

\begin{center}
\begin{tabular}{@{}lp{9cm}l@{}}
\toprule
Ergebnis & Aussage & Status \\
\midrule
Modul & \texttt{galois\_hrv\_preprocess.py} funktioniert & \KM \\
Bandgrenzen & 5-Limit-verträglich, nicht festgelegt & \BM \\
Nullmodell & Nahe Gittertreffer von LF/HF sind kein Beleg & \KM \\
BIDMC-01 & Negativ: zu kurz, Intensivpatient & \LM \\
Fantasia & Negativ: spontane Atmung, LF-CV 11--33\,\%, 4/10 im Zufallsbereich & \LM \\
Polar H10 & Peak auf der Atemfrequenz, HF = Harmonische, Treffer zufällig & \KM \\
Grund & LF/HF ist Leistungsverhältnis; Atmung bestimmt die Peaks & \BM \\
Nächstes & Prospektiver Kopplungstest, $n=30$, Blöcke zu 21~min & \SM \\
\bottomrule
\end{tabular}
\end{center}

\section*{Literatur}

\begin{description}[leftmargin=2em,labelwidth=1.8em]
\item[{[Pim18]}] Pimentel, M., Johnson, A., Charlton, P., \& Clifton, D.\ (2018).
  BIDMC PPG and Respiration Dataset (v1.0.0). PhysioNet.
  \url{https://doi.org/10.13026/C2208R}
\item[{[Iye96]}] Iyengar, N., Peng, C.\,K., Morin, R., Goldberger, A.\,L., \&
  Lipsitz, L.\,A.\ (1996). Age-related alterations in the fractal scaling of
  cardiac interbeat interval dynamics. \emph{Am.\ J.\ Physiol.\ Regul.\ Integr.\
  Comp.\ Physiol.}\ 271(4), R1078--R1084. Fantasia Database, PhysioNet,
  \url{https://doi.org/10.13026/C2RG61}
\item[{[Leh14]}] Lehrer, P.\,M., \& Gevirtz, R.\ (2014). Heart rate variability
  biofeedback: How and why does it work? \emph{Front.\ Psychol.}\ 5, 756.
\item[{[Sev22]}] Sevoz-Couche, C., \& Laborde, S.\ (2022). Heart rate
  variability and slow-paced breathing: When coherence meets resonance.
  \emph{Neurosci.\ Biobehav.\ Rev.}\ 135, 104576.
\item[{[Mak21]}] Makowski, D.\ et al.\ (2021). NeuroKit2: A Python toolbox for
  neurophysiological signal processing. \emph{Behav.\ Res.\ Methods} 53,
  1689--1696.
\item[{[Tas98]}] Tass, P.\ et al.\ (1998). Detection of $n{:}m$ phase locking
  from noisy data. \emph{Phys.\ Rev.\ Lett.}\ 81, 3291--3294.
\item[{[TF96]}] Task Force ESC/NASPE (1996). Heart rate variability.
  \emph{Eur.\ Heart J.}\ 17, 354--381.
\end{description}

\section*{Skripte}

\texttt{2/python/Dok360\_Skripte/}:
\texttt{galois\_hrv\_preprocess.py},
\texttt{fantasia\_galois\_analyse.py},
\texttt{orbit\_experiment\_synthetic.py},
\texttt{01\_rpeaks\_bidmc01.py},
\texttt{02\_welch\_5limit.py},
\texttt{03\_lombscargle\_dfa.py},
\texttt{pruef\_galois\_nullmodell.py}.
Polar-Auswertung, Nullmodell und alle Zahlen der Abschnitte~2 und~5:
\texttt{SWT\_Revision/pruef\_swt\_revision.py} (25/25 Assertions).
Fantasia-Daten: f1y01--05, f2o01--05 (.ecg/.hea, CC-BY~4.0).

\section*{Vermerk zur Verwendung von KI-Assistenz}

Bei der Erstellung dieses Dokuments wurde KI-Assistenz (Claude, Anthropic)
für Formulierung, LaTeX-Satz und Analyse-Skripte verwendet.
Die inhaltlichen Aussagen und die Verantwortung für den Text
liegen beim Autor.
